\section{Datasets and Sourcing}
\label{sec:data}

The analysis combines three categories of data: a primary set of eight
NDAP API tables published by the Reserve Bank of India, two
supplemental market datasets from Yahoo Finance, and a domain-curated
news corpus of seventy-five Indian monetary-policy events. Together
these cover all five channels through which overnight liquidity is
determined.

\subsection{Primary source: NDAP API (RBI open data)}

The eight NDAP datasets were fetched programmatically with
\texttt{scripts/stage1b\_fetch\_ndap.py}. The NDAP open-data API silently
paginates large responses, so the fetcher was built around an
auto-pagination loop that increments the \texttt{pageno} parameter and
appends responses until the API returns an empty page. Without this
behaviour the initial fetch would have truncated each series to the
first one or two hundred rows.

\begin{table}[H]
\caption{Primary datasets sourced from the NDAP API (RBI). Density is
  measured against the canonical Friday-aligned grid after Stage 2.}
\label{tab:ndap-sources}
\renewcommand{\arraystretch}{1.2}
\rowcolors{2}{SoftGray}{white}
\begin{tabularx}{\linewidth}{@{} l Y l c @{}}
\toprule
\rowcolor{Slate}\color{white}\textbf{Domain} &
\color{white}\textbf{NDAP table} &
\color{white}\textbf{Frequency} &
\color{white}\textbf{Cols retained} \\
\midrule
Monetary policy & RBI Weekly Statistics, Ratios and Rates & Weekly & 16 \\
Central-bank balance sheet & RBI Liabilities and Assets & Weekly & 22 \\
Money supply & RBI Weekly Aggregates (M3, Reserve Money) & Weekly & 3 \\
Debt markets & Commercial Paper Details & Weekly & 4 \\
Debt markets & Treasury Bills Details (91D, 182D, 364D) & Weekly & 12 \\
Repo markets & Market Repo Transactions & Weekly & 12 \\
Debt markets & Central Government Dated Securities & Weekly & 5 \\
Inflation & Major Price Indices (CPI) & Monthly $\to$ Weekly & 7 \\
\bottomrule
\end{tabularx}
\end{table}

The CPI table is the only monthly series. It is forward-filled to the
weekly grid with a maximum carry of four weeks, which is appropriate
because CPI is a stock variable that does not change between monthly
releases. Two CPI sub-indices were dropped after density filtering
(WPI at 63\% and one further sub-index at 37\%).

\subsection{External supplements: Yahoo Finance market data}

NDAP does not publish a machine-readable weekly equity index or
foreign-exchange series. Two supplemental sources were added through
the \texttt{yfinance} library in \texttt{scripts/stage1\_fetch\_yfinance.py}.

\begin{table}[H]
\caption{Yahoo Finance supplements. Both have 100\% density across the
analysis window after the indicator warmup period is trimmed.}
\label{tab:yfinance-sources}
\renewcommand{\arraystretch}{1.2}
\rowcolors{2}{SoftGray}{white}
\begin{tabularx}{\linewidth}{@{} l Y l l @{}}
\toprule
\rowcolor{Slate}\color{white}\textbf{Domain} &
\color{white}\textbf{Series} &
\color{white}\textbf{Ticker} &
\color{white}\textbf{Frequency} \\
\midrule
Equity markets & Nifty 50 OHLCV & \texttt{\textasciicircum NSEI} & Weekly \\
Foreign exchange & USD/INR OHLCV & \texttt{USDINR=X} & Weekly \\
\bottomrule
\end{tabularx}
\end{table}

The economic rationale for adding these series is that foreign
institutional equity flows directly affect rupee liquidity, and the
USD/INR level shapes the central bank's foreign-exchange intervention
decisions, both of which feed back into call-money supply. Including
these series also lets us empirically test whether market-sentiment
signals carry predictive content for overnight rates (hypothesis H4).

\subsection{Engineered features: technical indicators}

Five custom technical indicators are computed on top of each OHLCV
series in \texttt{scripts/stage2\_technical\_indicators.py}, producing ten
additional engineered columns per instrument and twenty-eight new
features in total.

\begin{itemize}
\item \textbf{Impulse MACD} (length 34, signal 9): a smoothed momentum
  oscillator with a histogram and signal line.
\item \textbf{SuperTrend} (period 7, multiplier 3.0): a trailing-stop
  trend filter encoded explicitly as $+1$, $-1$, or $0$ to preserve
  the warmup-period state.
\item \textbf{Squeeze Momentum} (Bollinger versus Keltner channel,
  conv 50, length 20): the squeeze correlation index.
\item \textbf{True Strength Index} (fast 13, slow 25, signal 13): a
  double-smoothed momentum oscillator with a signal line.
\item \textbf{Velocity} (lookback 14, EMA 9): rate-of-change indicator
  with smoothing.
\end{itemize}

These were sourced from \texttt{2013-01-01} to give 52 weeks of
warmup so that no leading NaN values bleed into the analysis window
that begins in February 2014.

\subsection{Curated NLP corpus of policy events}

The third data source is a hand-curated corpus of seventy-five Indian
monetary-policy events spanning February 2014 to August 2024,
constructed in \texttt{scripts/stage6\_news\_nlp.py}. Each event carries a
title, an English description, an event category (rate decision,
liquidity, structural, external, fiscal, or crisis), a sentiment score
in $[-1, +1]$, and an impact level. Source material was drawn from RBI
press releases, Bloomberg News, and Reuters coverage of MPC meetings
and inter-meeting actions.

\begin{warningbox}
This is a domain-expert labelled corpus, not a scraped news feed, and
is honestly framed as such throughout the report. The deliberate choice
to curate rather than scrape was driven by the need for clean, untyped
sentiment on a small number of high-impact events; an automated FinBERT
pipeline over five years of headlines is identified as future work in
Section~\ref{sec:limitations}.
\end{warningbox}

\subsection{Cross-domain combination, in line with NDAP guidance}

The project specification explicitly encourages combining datasets
across domains. The master panel does so by design: the eight NDAP
tables cover four distinct administrative domains (rates, balance
sheet, debt markets, inflation), and the Yahoo Finance and curated
news data add two more (capital markets and policy news). Together
these reflect the real-world mechanism by which WACMR is determined,
namely as the equilibrium outcome of five forces acting in parallel
rather than as a function of any single instrument.
